\documentclass[10pt,a4paper]{amsart}
\usepackage[margin=19mm]{geometry}
\usepackage{amsmath,amssymb,mathtools}
\usepackage[scr=rsfs,cal=euler]{mathalfa}
\usepackage{xfrac}
\usepackage[unicode,hidelinks,bookmarks=false]{hyperref}
\newcommand{\ie}{i.e.\ }
\newcommand{\cf}{cf.\ }
\newcommand{\resp}{respectively\ }
\newcommand{\Subsection}{\subsection}
\providecommand{\nobreakdash}{\nobreak\hbox{-}}
\setlength{\emergencystretch}{2em}
\multlinegap0pt
\usepackage{bm}
\usepackage{bbm}
\usepackage{dsfont}
\usepackage{xspace}
\usepackage{cancel}
\usepackage{mathtools}
\usepackage{amsthm}
\makeatletter
\@ifundefined{definition}{\newtheorem{definition}{Definition}}{}
\@ifundefined{lemma}{\newtheorem{lemma}[definition]{Lemma}}{}
\@ifundefined{remark}{\newtheorem{remark}[definition]{Remark}}{}
\@ifundefined{theorem}{\newtheorem{theorem}[definition]{Theorem}}{}
\makeatother
\def\dxy{\,dx dy}
\def\dx{\,dx}
\def\e{\varepsilon}
\title[Nonlocal-interaction vortices]{Nonlocal-interaction vortices}
\author{Margherita Solci}
\date{Source: arXiv:2302.06526v2 (CC BY 4.0)}
\begin{document}
\maketitle

\noindent\emph{Source and licence.} Nonlocal-interaction vortices. Version arXiv:2302.06526v2 (CC BY 4.0), distributed under the Creative Commons Attribution 4.0 International licence, as recorded in the arXiv licence metadata for this exact version and shown on the abstract page https://arxiv.org/abs/2302.06526v2. Full rights notice: licenses/arxiv-2302.06526-CC-BY-4.0.txt.\par\smallskip

\noindent\textit{Editorial scope.} Counted unit: the Definition whose source label is \texttt{conv} in the article (source0.tex lines 303--310, source label \texttt{conv}), delivered together with the article's own complete original proof of it (source0.tex lines 619--764, one \texttt{proof} environment of 146 lines). No mathematics is written, repaired or re-derived: statement and proof are the article's own text line for line. The proof is detached from the statement in the source: the article states the result at lines 303--310 and prints its proof at lines 619--764, and the proof environment's own header names the statement, so the association is the article's own. Required context retained uncounted: lemma2 (lines 252--261); def-ae (lines 294--296); comp-convergence (lines 322--348); effe-e-d-comp (lines 362--364); compactness (lines 367--374); e1 (lines 398--404); conf-Aeu-g (lines 478--481); stima-z (lines 490--500); conf-Aeu (lines 510--515); indepp (lines 546--600); uexiz (lines 569--571); defxi (lines 587--587). Every definition, display and label of those lines is retained unchanged. Omitted from this package and not redistributed: 1--251, 262--293, 297--302, 311--321, 349--361, 365--366, 375--397, 405--477, 482--489, 501--509, 516--545, 572--586, 588--618, 765--1021. Nothing inside a retained excerpt is omitted or altered: every retained body is delivered exactly as the article's lines are written, so no omission receipt of any kind accompanies this package. Citations: the retained excerpts carry no citation command at all, so no bibliography entry of the article is transcribed and no reference is redistributed. Reference adaptation: none was needed and none was made; every reference inside the delivered unit resolves inside it, so no unresolved reference or citation is produced. Printed slips kept literal: no printed word, symbol, hypothesis or number of the counted statement or of its proof was altered. Layout and class: the article's own class is \texttt{article} with packages including geometry, amsmath, amsthm, amssymb, graphicx, color, epstopdf, pdfsync, mathabx; the wrapper is the lane's pinned \texttt{amsart} editorial wrapper at 10pt on A4, which restates the theorem-like environments the retained text needs and adds the article's own macro definitions that the retained text needs (\texttt{\textbackslash dx}, \texttt{\textbackslash dxy}, \texttt{\textbackslash e}), with extra wrapper packages (bm, bbm, dsfont, xspace, cancel, mathtools). The delivered numbering is the unit's own. Assets: no retained span contains a figure environment, a graphics-inclusion command, a PDF-page inclusion, a table or a TikZ picture; no image file is copied into this package, the package declares no asset dependency and no image file is redistributed with it. Licence: version arXiv:2302.06526v2 of this article is distributed under the Creative Commons Attribution 4.0 International licence, as recorded in the arXiv licence metadata for this exact version and shown on the abstract page https://arxiv.org/abs/2302.06526v2. A full rights notice is delivered at licenses/arxiv-2302.06526-CC-BY-4.0.txt. Characters: every retained excerpt is pure ASCII, so the delivered engine, XeLaTeX with its default Latin Modern text font, reports no missing character.
\begin{lemma}[coerciveness of discrete energies] \label{lemma2}
Let $\{u^\e\}$ be a family of discrete functions such that $\sup_\e X_\e(u^\e)<+\infty$; then, there exist a subsequence $\e_j$ and an integral $(d-2)$-current $M$   
such that 
\begin{equation}\label{codico}
\lim_{j\to +\infty}{\bf F}_U(\star JA_{\e_j}(u^{\e_j})-\mu)=0 
\end{equation}
for all $U\subset\subset\Omega$ open sets, where $\mu=\pi M$. 
In particular, if $d=2$ there exist $N\in \mathbb N$, $x_\ell\in \Omega$, and $d_\ell\in\mathbb Z$ for $\ell\in\{1,\ldots,N\}$ such that $\mu$ is an atomic measure given by 
$$\mu=\pi\sum_{\ell=1}^N d_\ell \delta_{x_\ell}.$$ 
\end{lemma}

\begin{equation}\label{def-ae}
A_\e(v)(x)=\sum_{k\in\mathbb Z^d}\lambda_k\Big({x\over\e}\Big) v(\e k). 
\end{equation} 

\begin{theorem}[Compactness and $\Gamma$-convergence]\label{comp-convergence}
Let $\rho\colon[0,+\infty)\to[0,+\infty)$ be a fixed kernel with compact support such that 
\begin{enumerate}
    \item[\rm(a)] $\displaystyle\int_{\mathbb R^d}\rho(|\xi|)|\xi|^2\, dx<+\infty;$
    \item[\rm(b)] $\rho(t)\geq \rho_0 >0$ in a neighbourhood of $0$. 
\end{enumerate} 
Let $\Omega\subset\mathbb R^d$ be an open Lipschitz bounded domain, and for all $\e>0$ let $F_\e\colon L^1(\Omega; S^1)\to [0,+\infty]$ be defined by 
\begin{equation*}
F_\e(u)=
{1\over\e^{d+2}|\log\e|} \int_{\Omega\times\Omega}\rho\Bigl({|x-y|\over\e}\Bigr) {|u(x)-u(y)|^2}\dxy.
\end{equation*}
Then, we have the following results. 
\begin{enumerate}
\item[{\rm (i)}] {\rm (equi-coerciveness of $F_\e$)} 
If $\{u_\e\}$ is a sequence such that 
$F_\e(u_\e)$ is equibounded,  
then, up to subsequences, there exists an integral $(d-2)$-current $M$ 
such that $\{u_\e\}$ converges to $\mu=\pi M$ in the sense of Definition {\rm\ref{conv}}. 
\item[{\rm (ii)}] {\rm (lower bound)} 
If $\{u_\e\}$ converges to $\mu=\pi M$ in the sense of Definition {\rm\ref{conv}}, then 
\begin{equation}\label{Crho}
\liminf_{\e\to 0}F_\e(u_\e)\geq C_\rho \|M\|, \ \ \hbox{ where } \  C_\rho={2\pi\over d}\int_{\mathbb R^d}\rho(|\xi|)|\xi|^2\, d\xi.
\end{equation} 
\item[{\rm (iii)}] {\rm (upper bound)} For every integral $(d-2)$-current $\mu=\pi M$, there exists a sequence $\{u_\e\}$ converging to $\mu=\pi M$ in the sense of Definition {\rm\ref{conv}} such that 
$$\lim_{\e\to 0}F_\e(u_\e)= C_\rho \|M\|.$$
\end{enumerate}
\end{theorem}

\begin{equation}\label{effe-e-d-comp}
F_\e(u)={1\over\e^{d+2}|\log\e|} \int_{\Omega\times\Omega}\chi_{[-1,1]^d}\Bigl({x-y\over\e}\Bigr) {|u(x)-u(y)|^2}\dxy 
\end{equation}

\begin{lemma}[Compactness]\label{compactness}
Let $\Omega\subset \mathbb R^d$ be a bounded open Lipschitz domain and let 
$\{u_\e\}$ be a sequence such that 
$F_\e(u_\e)$ is equibounded, where $F_\e$ is defined by
\eqref{effe-e-d-comp}. 
Then, up to subsequences, there exists an integral $(d-2)$-current $M$ 
such that $\{u_\e\}$ converges to $\mu=\pi M$ in the sense of Definition {\rm\ref{conv}}. 
\end{lemma}

\begin{eqnarray}\label{e1}
&&\hspace{-5mm}{1\over\e^2|\log\e|}\int_{U} |u(x+\e e_1)-u(x)|^2\dx\nonumber \\
&&\leq 
{2\over\e^2|\log\e|}\int_{U} \Big|u(x+\e e_1)-u\Big(x+\frac{\e}{2} e_1\Big)\Big|^2\dx+
{2\over\e^2|\log\e|}\int_{U} \Big|u\Big(x+\frac{\e}{2} e_1\Big)-u(x)\Big|^2\dx\nonumber \\
&&\leq C\, F_\e(u)
\end{eqnarray} 

\begin{equation}\label{conf-Aeu-g}
\frac{1}{|\log\e|}
\int_{U} |\nabla A_\e(I_\e(u_\e))(x)|^2\, dx\le C F_\e(u_\e). 
\end{equation} 

\begin{eqnarray}\label{stima-z}
&&\hspace{-1cm}\int_{[0,1]^2}\int_{U}|u_\e^z(x)-u_\e(x)|^2\, dx\, dz\nonumber \\
&&=\int_{[0,1]^2}\int_{U}\Big|\sum_{k\in\mathbb Z^2}\lambda_k\Big(\frac{x-\e z}{\e}\Big) \big(u_\e(\e (k+z))-u_\e(x)\big)\Big|^2\, dx\, dz\nonumber\\
&&\leq \int_{[0,1]^2}\int_{U}\sum_{k\in\mathbb Z^2}\lambda_k\Big(\frac{x-\e z}{\e}\Big) \Big|u_\e(\e (k+z))-u_\e(x)\Big|^2\, dx\, dz\nonumber\\
&&=\int_{U} \int_{[0,1]^2}\sum_{k\in\mathbb Z^2}\lambda_k\Big(\frac{x-\e z}{\e}\Big) \Big|u_\e(\e (k+z))-u_\e(x)\Big|^2\, dz\, dx\nonumber\\
&&\leq\sum_{j\in\mathcal I_\e}\int_{Q_\e^j} \int_{[0,1]^2}\sum_{k\in\mathbb Z^2}\lambda_k\Big(\frac{x-\e z}{\e}\Big) \Big|u_\e(\e (k+z))-u_\e(x)\Big|^2\, dz\, dx\nonumber\\
&&\leq\sum_{j\in\mathcal I_\e}\int_{Q_\e^j} \frac{C}{\e^2}\int_{x+[-\e,\e]^2} |u_\e(y)-u_\e(x)|^2\, dy\, dx\nonumber\\
&&\leq\frac{C}{\e^2}\int_{\Omega\times\Omega}\chi_{[-1,1]^2}\Big(\frac{x-y}{\e}\Big) |u_\e(y)-u_\e(x)|^2\, dy\, dx\nonumber\\
&& 
=C\e^2|\log\e| F_\e(u_\e). 
\end{eqnarray} 

\begin{eqnarray}\label{conf-Aeu}\nonumber
\int_U| A_\e(I_\e(u_\e))(x) - u_\e(x)|^2\, dx
&=&\int_U\Big|\int_{[0,1]^2} (u_\e^z (x+\e z) - u_\e(x))\, dz\Big|^2\, dx\\\nonumber
&\leq&\int_U\int_{[0,1]^2} |u_\e^z (x+\e z) - u_\e(x)|^2\, dz\, dx\\
&\leq& C\e^2|\log \e| F_\e(u_\e).
\end{eqnarray}

\begin{remark}[Independence from the discretization]\label{indepp}\rm 
We now remark that Definition \ref{conv} is in fact independent from the choice of the discretization. 
We make this statement precise in the case $d=2$, the general case following with minor modifications. 

For a fixed $\xi=(\xi_1,\xi_2)\in\mathbb R^2\setminus\{0\}$
we consider the lattice $\mathbb Z_\xi^2=\mathbb Z \xi \oplus \mathbb Z \xi^\perp$, 
where $\xi^\perp=(\xi_2,-\xi_1)$ and the convergence of a sequence $\{u_\e\}$ obtained by a discretization on this lattice.   
Namely, for each $u\in L^1_{\rm loc}(\Omega; \mathbb R^2)$ and $\e>0$ we define $I^\xi_\e(u)\colon\e\mathbb Z^2_\xi \to\mathbb R^2$ as follows
\begin{equation}
I_\e^\xi(u)(\e k)=
\frac{1}{|\xi|^2\e^2}\int_{Q_\e^{k,\xi}\cap \Omega}u(x)\, dx,
\end{equation}
where $Q_\e^{k,\xi}=\e k+ \e Q^\xi$, $k\in \mathbb Z^2_\xi$ and $Q^\xi=[0,1]\xi+[0,1]\xi^\perp$.    

Then, we consider a fixed triangulation of $\mathbb R^2$ with vertices in $\mathbb Z^2_\xi$ and the corresponding family of piecewise-affine functions $\{\lambda_k^\xi\}_{k\in\mathbb Z^2_\xi}$ with $\lambda_k^\xi\colon\mathbb R^2\to[0,1]$ and such that $\lambda_k^\xi(k)=1$, the support of $\lambda_k^\xi$ is the union of the elements of the triangulation containing $k$ and $\sum_{k\in\mathbb Z^2_\xi}\lambda_k(x)=1$ for all $x$. 
Given a discrete function $v\colon\e \mathbb Z^2_\xi\to\mathbb R^2$, following \eqref{def-ae} we define a piecewise-affine interpolation by setting
\begin{equation}
A_\e^\xi(v)(x)=\sum_{k\in\mathbb Z^2_\xi}\lambda_k^\xi\Big({x\over\e}\Big) v(\e k). 
\end{equation}

The key argument of Lemma \ref{compactness}  is the comparison of the functions $u_\e^z$ with $u_\e$ as in \eqref{stima-z}, and is obtained thanks to \eqref{e1}. 
We can repeat the arguments leading to \eqref{e1} with $\xi$ in the place of $e_1$, up to changing the constants, as follows. 
For fixed $\xi$ and $z\in[0,1]^2$, for $i\in\mathbb Z^2$ we set 
\begin{equation}\label{uexiz}
u^{\e,z,\xi}_i=u^{\e,z,\xi}(\e i)=u_\e(\e z_1\xi+\e z_2\xi^\perp+\e i_1 \xi +\e i_2 \xi^\perp), 
\end{equation} 
and let $u^{z,\xi}_\e$ denote the corresponding piecewise-affine interpolation from the lattice $\e z_1\xi+\e z_2\xi^\perp+\e \mathbb Z^2_\xi$. 
Note that $u^{\e,z,\xi}_i$ are defined for $i\in\mathbb Z^2$, so that they differ form $u_\e^{z,\xi}$ by the linear transformation carrying $\mathbb Z^2$ in $\mathbb Z^2_\xi$, up to a small translation.  
Then we obtain the estimates analogous to \eqref{conf-Aeu-g} and \eqref{conf-Aeu} with $A^\xi_{\e}(I^\xi_{\e}(u_{\e}))$ in the place of $A_{\e}(I_{\e}(u_{\e}))$, and we can conclude that 
\begin{eqnarray}\nonumber&&\hskip-1.5cm\int_{[0,1]^2} X_{\e} (u^{\e,z,\xi};U^\e_{z,\xi})\,dz +\frac{1}{\e^2|\log\e|}\int_{[0,1]^2} \int_U| A^\xi_{\e}(I^\xi_{\e}(u_{\e} (x))) - u^{z,\xi}_{\e}(x)|^2\, dx\,dz\\
&&\hskip1.5cm+
 \frac{1}{|\log{\e}|}\int_{[0,1]^2} \int_U| \nabla A^\xi_{\e}(I^\xi_{\e}(u_{\e})) - \nabla u^{z,\xi}_{\e}(x)|^2\, dx\,dz\le C(|\xi|) F_\e(u_\e), 
 \end{eqnarray}
where
$U^\e_{z,\xi}=\{x\in\mathbb R^2: \e z_1 \xi+\e z_2\xi^\perp +L(\xi)x\in U\}$ 
and $L(\xi)$ is the linear map such that $L(\xi)(e_1)=\xi$ and $L(\xi)(e_2)=\xi^\perp$. 

As a result, concluding as in the proof of the proposition, we obtain 
a subsequence $\e_j=\e_j(\xi)$ and 
a measure 
$\mu(\xi)$ such that 
\begin{equation}\label{defxi}\lim_{j\to+\infty}{\bf F}_U(\star J(A^\xi_{\e_j}(I_{\e_j}^\xi(u_{\e_j})))-\mu(\xi))=0\end{equation} 
for all $U\subset\subset\Omega$. 
Note that if $\mu(\xi)=\pi\sum_{\ell=1}^Nd^\xi_\ell\delta_{x^\xi_\ell}$, then $u^{\e_j,z,\xi}$ converges to 
$\pi\sum_{\ell=1}^Nd^\xi_\ell\delta_{L^{-1}(x^\xi_\ell)}$. 

Now, we can apply Lemma \ref{lemma2} to $A^\xi_{\e_j}(I_{\e_j}^\xi(u_{\e_j}))$ and $A_{\e_j}(I_{\e_j}(u_{\e_j}))$ with the estimate on the gradients given by \eqref{conf-Aeu-g}, while the second assumption can be obtained from \eqref{conf-Aeu} by a triangular argument. 
Hence, the sequence $\{u_\e\}$ converges up to subsequences to the measure $\mu(\xi)$ in the sense of Definition \ref{conv},  
showing that the convergence indeed depends only on $u_\e$ and not on the discretization chosen.
Indeed, if the sequence $u_\e$ converges to $\mu$ in the sense of Definition \ref{conv}, 
then $u_{\e_j}$ converges to $\mu$, implying $\mu=\mu(\xi)$. 
This allows to remark that, if in Definition \ref{conv} we require that \eqref{defxi} holds 
for all $U\subset\subset\Omega$, 
we obtain a definition of convergence which is equivalent to Definition \ref{conv}. 
\end{remark}

\begin{definition}[Convergence]\label{conv}
Let $\{u_\e\}$ be a sequence in $L^2_{\rm loc}(\Omega;\mathbb R^2)$ and let $\mu=\pi M$, where 
$M$ is an integral $(d-2)$-current.  
The sequence $\{u_\e\}$ {\rm converges to} $\mu$ if 
for every $U\subset\subset\Omega$ we have  
$$\lim_{\e\to 0}{\bf F}_U(\star J(A_\e(I_\e(u_\e)))-\mu)=0,$$
where ${\bf F}_U$ is the flat norm in $U$ and $\star J$ is the Jacobian current. 
\end{definition}

\begin{proof}[Proof of Theorem {\rm\ref{comp-convergence}} {\rm (ii)} (lower bound)] 
In order to point out and clarify the key steps of the proof, we first deal with the case $d=2$. 
\smallskip 

\noindent(a) {\em The case $d=2$.} \ In this case, we can consider an integral $0$-current of the form 
$M=\sum_{\ell=1}^nd_{\ell}\delta_{x_\ell}$. 
Let $u_\e$ converge to $\mu=\pi M$ in the sense of Definition 
\ref{conv}. 
Let $U\subset\subset\Omega$. In analogy with the definition of the family of indices $\mathcal I_\e$, we set 
$$\mathcal I_\e^\xi(U)=\mathcal I_\e^\xi=\{k\in\mathbb Z^2_\xi: \e k+4\e \widehat Q^\xi\subset\subset \Omega, \ 
\e k+2\e \widehat Q^\xi\cap U\neq\emptyset\},$$ 
where $\widehat Q^\xi$ is the square centered at $0$ with edges $\xi$ and $\xi^\perp$. 
We have 
\begin{eqnarray}\label{primainf}
F_\e(u_\e)&=&\int_{\mathbb R^2}\rho(|\xi|) \Bigl({1\over\e^2|\log\e|}\int_{\Omega} |u_\e(x+\e\xi)-u_\e(x)|^2\dx\Bigr)\,d\xi\nonumber\\
&\geq& \frac{1}{2|\log\e|}\int_{\mathbb R^2}\rho(|\xi|)|\xi|^2 
\Bigl(\int_{[0,1]^2}\sum_{k\in\mathcal I_\e^\xi} |u_\e(\e s_1 \xi+ \e s_2\xi^\perp+\e k+\e \xi)\nonumber\\
&&\hspace{5cm}-u_\e(\e s_1 \xi+\e s_2\xi^\perp +\e k)|^2 \, ds\Bigr)\,d\xi\nonumber\\
&&+\frac{1}{2|\log\e|}\int_{\mathbb R^2}\rho(|\xi|)|\xi|^2 
\Bigl(\int_{[0,1]^2}\sum_{k\in\mathcal I_\e^\xi} |u_\e(\e s_1 \xi+\e s_2\xi^\perp+\e k+\e \xi^\perp)\nonumber\\
&&\hspace{5cm}-u_\e(\e s_1 \xi+\e s_2\xi^\perp +\e k)|   \, ds\Bigr)\,d\xi\nonumber\\
&\geq& \frac{1}{2|\log\e|}\int_{\mathbb R^2}\rho(|\xi|)|\xi|^2 
\Bigl(\int_{[0,1]^2}\frac{1}{2}\sum_{\langle i, j\rangle} |u_i^{\e,s,\xi}-u_j^{\e,s,\xi}|^2 \, ds\Bigr)\,d\xi\nonumber\\
&\geq& \frac{1}{4}\int_{\mathbb R^2}\rho(|\xi|)|\xi|^2 
\Bigl(\int_{[0,1]^2}X_\e (u^{\e,s,\xi}; U^\e_{s,\xi}) \, ds\Bigr)\,d\xi
\end{eqnarray}
where for $i\in\mathbb Z^2$ the value $u^{\e,s,\xi}_i$ is defined as in \eqref{uexiz} and  
$$U^\e_{s,\xi}=\{x\in\mathbb R^2: \e s_1 \xi+\e s_2\xi^\perp +L(\xi)(x)\in U\},$$  
 $L(\xi)$ being the linear map such that $L(\xi)(e_1)=\xi$ and $L(\xi)(e_2)=\xi^\perp$. 

We now introduce a parameter $\eta>0$ and write
\begin{eqnarray*}
F_\e(u_\e)& =&
(1-\eta)F_\e(u_\e) +\eta F_\e(u_\e)\\
&\ge& \frac{1}{2}\int_{\mathbb R^2\setminus B_\eta}\rho(|\xi|)|\xi|^2 
\int_{[0,1]^2}(1-\eta)\Bigl(X_\e (u^{\e,s,\xi};U_\xi) \\
&&\hskip3cm +C_\eta\frac{1}{\e^2|\log\e|}\int_U| A_\e(I_\e(u_\e (x))) - u^{s,\xi}_\e(x)|^2\, dx\nonumber\\
 &&\hskip3cm +C_\eta
 \frac{1}{|\log\e|}\int_U| \nabla A_\e(I_\e(u_\e)) - \nabla u^{s,\xi}_\e(x)|^2\, dx\Big)\, ds\,d\xi,
\end{eqnarray*}
where $U_\xi$ is any open set contained in the intersection of the sets $U^\e_{s,\xi}$ for $s\in[0,1]^2$ and $\e$ small enough, $u^{s,\xi}_\e$ denotes the corresponding piecewise-affine interpolation from the lattice $\e s_1\xi+\e s_2\xi^\perp+\e \mathbb Z^2_\xi$ as in Remark \ref{indepp} and $C_\eta$ is a positive constant.

By applying Fatou's Lemma we note that for almost all $\xi$ and $s$ the limit 
\begin{eqnarray*}
&&L(\xi,s)=\liminf_{\e\to 0} \Bigl(X_\e (u^{\e,s,\xi};U_\xi) +C_\eta\frac{1}{\e^2|\log\e|}\int_U| A_\e(I_\e(u_\e (x))) - u^{s,\xi}_\e(x)|^2\, dx\nonumber\\
 &&\hskip5cm +C_\eta
 \frac{1}{|\log\e|}\int_U| \nabla A_\e(I_\e(u_\e)) - \nabla u^{s,\xi}_\e(x)|^2\, dx\Big)
\end{eqnarray*}
is finite. Hence we can find a sequence $\e_j$ (depending on $\xi$ and $s$) such that 
\begin{eqnarray*}
&&L(\xi,s)=\lim_{j\to+\infty} \Bigl(X_{\e_j} (u^{\e_j,s,\xi};U_\xi)+C_\eta\frac{1}{\e_j^2|\log\e_j|}\int_U| A_{\e_j}(I_{\e_j}(u_{\e_j} (x))) - u^{s,\xi}_{\e_j}(x)|^2\, dx\\&&\hskip5cm +C_\eta
 \frac{1}{|\log{\e_j}|}\int_U| \nabla A_{\e_j}(I_{\e_j}(u_{\e_j})) - \nabla u^{s,\xi}_{\e_j}(x)|^2\, dx\Big)
\end{eqnarray*}
and equals the liminf above.  
By Remark \ref{indepp} we deduce that $u^{\e_j,s,\xi}$ converges to 
$$\hat\mu(\xi)=\pi\sum_{\ell=1}^Nd_\ell\delta_{L^{-1}(\xi)(x_\ell)},$$ 
so that
$$L(\xi,s)\ge \liminf_{j\to+\infty}X_{\e_j} (u^{\e_j,s,\xi};U_\xi)\ge 4\pi\sum_{\{\ell: x_\ell\in U\}} |d_\ell|$$ 
since we may assume that $L^{-1}(\xi)(x_\ell)\in U_\xi$. 

We can then proceed in the application of Fatou's Lemma to deduce that
$$\liminf_{\e\to0}F_\e(u_\e)\ge (1-\eta)\int_{\mathbb R^2\setminus B_\eta}\rho(|\xi|)|\xi|^2d\xi\,  \pi\sum_{\{\ell: x_\ell\in U\}}|d_\ell|,
$$
and finally, using the arbitrariness of $\eta$ and $U$, that
$$\liminf_{\e\to0}F_\e(u_\e)\ge \int_{\mathbb R^2}\rho(|\xi|)|\xi|^2d\xi\,  \pi\sum_{\ell=1}^N|d_\ell|,$$
which is the desired lower bound.

\smallskip 

\noindent(b) {\em The general case.}\ 
In the $d$-dimensional case we cannot simply take into account orthogonal bases of the form $\{\xi, \xi^\perp\}$. 
In order to repeat the argument in the $2$-dimensional case using the lower estimate for the discrete functionals $X_\e$, we will consider the space of the orthonormal bases in $\mathbb R^d$ 
and use them to parameterize the interpolations, generalizing the role of $\xi$ and $\xi^\perp$ in the $2$-dimensional computations.  
To this end, we define the space
$$V=\{\overline \nu=(\nu_1,\dots, \nu_d): \nu_j\in S^{d-1}  \ \hbox{\rm such that }  \langle \nu_i, \nu_j\rangle=0 \ \hbox{\rm for } i\neq j\},$$ 
whose Hausdorff dimension $\frac{d(d-1)}{2}$ is denoted by $k_d$. 
Moreover, for any $\nu\in S^{d-1}$ and $n\in \{1,\dots, d\}$ we set 
$$V^{\nu}_n=\{\overline \nu=(\nu_1,\dots, \nu_d)\in V: \nu_n=\nu\}.$$ 
The Hausdorff dimension of $V^{\nu}_n$ is $k_d-(d-1)$ and we have 
\begin{equation}\label{dim} 
\mathcal H^{k_d-(d-1)}(V^{\nu}_n)=\frac{\mathcal H^{k_d}(V)}{\mathcal H^{d-1}(S^{d-1})}.
\end{equation} 
Note that in fact the space $V$ is the orthogonal group $O(d)\subset GL(d)$; that is, the group of the orthogonal $d\times d$ matrices, and each $V^{\nu}_n$ corresponds to the orthogonal group $O(d-1)$ acting on the orthogonal complement of $\nu$. 

Now, let $M$ be an integral $(d-2)$-current, and let the sequence $\{u_\e\}$ converge to $\mu=\pi M$ in the sense of Definition \ref{conv}. 
Let $U\subset\subset\Omega$. In analogy with the previous case, for $t>0$ and $\overline \nu\in V$ we set 
$$\mathcal I_\e^{t\overline \nu}(U)=\mathcal I_\e^{t\overline \nu}=\{k\in\mathbb Z^d_{t\overline \nu}: \e k+4\e \widehat Q^{t\overline \nu}\subset\subset \Omega, \ 
\e k+2\e \widehat Q^{t\overline \nu}\cap U\neq\emptyset\},$$ 
where $\mathbb Z^d_{t\overline \nu}=
\mathbb Z t\nu_1 \oplus \dots \oplus \mathbb Z t\nu_1$ and $\widehat Q^{t\overline \nu}$ is the square centered at $0$ with edges $\{t\nu_1\}_{n=1}^{d}$. Again note that $\{t\nu_1, \dots, t\nu_d\}$ plays the same role as $\{\xi, \xi^\perp\}$ in the $2$-dimensional case. 
As in \eqref{primainf}, by using \eqref{dim} we have 
\begin{eqnarray}\label{stimainfd}
F_\e(u_\e)&=&\int_{\mathbb R^d}\rho(|\xi|) \Bigl({1\over\e^2|\log\e|}\int_{\Omega} |u_\e(x+\e\xi)-u_\e(x)|^2\dx\Bigr)\,d\xi\nonumber\\
&=&\int_0^{+\infty} \rho(t) t^{d-1}\int_{S^{d-1}} \frac{1}{\e^2|\log\e|} 
\int_{\Omega} |u_\e(x+\e t\nu)-u_\e(x)|^2\dx \, d\mathcal H^{d-1}(\nu)\, dt
\nonumber\\
&=&\int_0^{+\infty} \rho(t) t^{d-1}\int_{S^{d-1}} \frac{1}{d}\sum_{n=1}^d \frac{\mathcal H^{d-1}(S^{d-1})}{\mathcal H^{k_d}(V)}\int_{V^{\nu}_n} \frac{1}{\e^2|\log\e|}  \nonumber 
\\ 
&&{\hspace{1cm}}
\int_{\Omega} |u_\e(x+\e t\nu)-u_\e(x)|^2\dx \, d\mathcal H^{k_d-(d-1)}(\overline \nu_n)\, d\mathcal H^{d-1}(\nu)\, dt\nonumber \\
&=&\frac{1}{d}\frac{\mathcal H^{d-1}(S^{d-1})}{\mathcal H^{k_d}(V)} \int_0^{+\infty} \rho(t) t^{d-1} \int_{V} \frac{1}{\e^2|\log\e|}  
\nonumber 
\\ 
&&{\hspace{1cm}}
\int_{\Omega}\sum_{n=1}^d |u_\e(x+\e t\nu_n)-u_\e(x)|^2\dx \, d\mathcal H^{k_d}(\overline \nu)\, dt,  
\end{eqnarray}
where $\overline \nu_n=(\nu_1, \dots, \nu_{n-1},\nu_{n+1},\dots, \nu_d)\in (S^{d-1})^{d-1}$ parameterizes an element of 
$V^{\nu}_n$ by not considering the $n$-th component $\nu$. 
Now, proceeding as in the $2$-dimensional case, we subdivide $\Omega$ in cubes corresponding to indices in $\mathcal I_\e^{t\overline \nu}$, obtaining 
\begin{eqnarray*}
&&\hspace{-5mm}\frac{1}{\e^2|\log\e|}  \int_{\Omega}\sum_{n=1}^d |u_\e(x+\e t\nu_n)-u_\e(x)|^2\dx \\
&&\geq t^d \int_{[0,1]^d} \frac{1}{|\log\e|}\sum_{k\in \mathcal I_\e^{t\overline \nu}} \e^{d-2} 
\Big|u_\e\Big(\e \sum_{\ell=1}^d t s_\ell\nu_\ell +\e t\nu_n +\e k\Big)-u_\e\Big(\e \sum_{\ell=1}^d t s_\ell\nu_\ell +\e k\Big)\Big|^2\, ds\\
&&=t^d \int_{[0,1]^d} \frac{1}{|\log\e|} \frac{1}{2}\sum_{\langle i, j \rangle} |u_i^{\e, s, t\overline \nu}-u_j^{\e, s, t\overline \nu}|^2 \, ds\\
&&\geq \frac{t^d}{2} \int_{[0,1]^d} X_\e(u^{\e,s, t\overline \nu}; U^\e_{s,t\overline\nu}) \, ds, 
\end{eqnarray*} 
where, for fixed $t>0$, $\overline \nu\in V$ $s\in[0,1]^d$ and $i\in\mathbb Z^d$, the value $u^{\e,s,t\overline \nu}_i$ is defined 
by
\begin{equation}\label{uexizd}
u^{\e,s,t\overline \nu}_i=u^{\e,s,t\overline \nu}(\e i)=u_\e\bigg(\e\sum_{\ell=1}^d t s_\ell\nu_\ell 
+\e \sum_{\ell=1}^d t i_\ell \nu_\ell\bigg)
\end{equation} 
and the set $U^\e_{s,t\overline \nu}$ is given by 
$$U^\e_{s,t\overline \nu}=\Big\{x\in\mathbb R^d: \e\sum_{\ell=1}^d t s_\ell \nu_\ell+t L(\overline u)(x)\in U\Big\},$$ 
with $L(\overline \nu)$ the linear map such that $L(\overline \nu)(e_\ell)=\nu_\ell$ for any $
\ell\in\{1,\dots, d\}$. 
Hence, from \eqref{stimainfd} we get the estimate 
\begin{equation}\label{stimainfd2}
F_\e(u_\e)\geq \frac{1}{2d}\frac{\mathcal H^{d-1}(S^{d-1})}{\mathcal H^{k_d}(V)} 
\int_{V}\int_0^{+\infty} \rho(t) t^{2d-1} \int_{[0,1]^d} X_\e(u^{\e,s, t\overline \nu}; U^\e_{s,t\overline\nu}) \, ds   \, dt\, d\mathcal H^{k_d}(\overline \nu).  
\end{equation}
Note moreover that, as in Remark \ref{indepp}, the functions $u^{\e,s,t \overline \nu}$ converge to the pull-back of the limit measure $\mu$ with respect to $t L(\overline \nu)$, so that 
$$\liminf_{\e\to 0} X_\e (u^{\e,s,t\overline \nu}; U^\e_{s,t\overline \nu})\geq 4\pi t^{2-d}\|M\|.$$


Now we can conclude the proof of the lower inequality, obtaining 
\begin{eqnarray*}
\liminf_{\e\to 0}F_\e(u_\e)&\geq& 
\frac{2\pi}{d}\frac{1}{\mathcal H^{k_d}(V)} 
\int_{V}
\int_0^{+\infty} \mathcal H^{d-1}(t S^{d-1}) \rho(t) t^{2}    \, dt\, d\mathcal H^{k_d}(\overline \nu)\ \|M\|(U)
\\
&\geq&  \frac{2\pi}{d}\int_{\mathbb R^d}\rho(|\xi|)|\xi|^2d\xi\,  \|M\|(U),
\end{eqnarray*} 
by exactly following the steps in the $2$-dimensional case. 
\end{proof}

\end{document}
